\documentclass[11pt,a4paper]{article}
\usepackage[italian]{babel}
\usepackage[margin=2.5cm]{geometry}
\usepackage{parskip}
\usepackage{amsmath, amssymb, amsthm}
\usepackage{booktabs}
\usepackage{array}
\usepackage{xcolor}
\usepackage{needspace}
\usepackage{enumitem}
\usepackage{listings}
\usepackage{tikz}
\usepackage[most]{tcolorbox}
\usepackage[hidelinks]{hyperref}
\usetikzlibrary{decorations.pathreplacing, shapes.geometric, arrows.meta, positioning}

% Niente vedove/orfane: i paragrafi non lasciano righe isolate a fine/inizio pagina
\widowpenalty=10000
\clubpenalty=10000

% Elenchi più compatti
\setlist{itemsep=2pt, parsep=0pt, topsep=4pt, partopsep=0pt}

% --- Riquadri con bordo nero, senza numero (si spezzano tra due pagine) ---
% Uso: \begin{definizione} ... \end{definizione}
%      \begin{teorema}[Nome] ... \end{teorema}  (il nome è facoltativo)
\tcbset{nota/.style={enhanced jigsaw, breakable, colback=white,
  colframe=black, boxrule=0.5pt, sharp corners}}
\NewTColorBox{definizione}{}{nota,
  before upper={\textbf{Definizione.}\ }}
\NewTColorBox{teorema}{o}{nota,
  before upper={\textbf{Teorema\IfValueT{#1}{ (#1)}.}\ }}
\NewTColorBox{proposizione}{o}{nota,
  before upper={\textbf{Proposizione\IfValueT{#1}{ (#1)}.}\ }}
\NewTColorBox{proprieta}{o}{nota,
  before upper={\textbf{Proprietà\IfValueT{#1}{ (#1)}.}\ }}
\NewTColorBox{assioma}{o}{nota, boxrule=1pt,
  before upper={\textbf{Assioma\IfValueT{#1}{ (#1)}.}\ }}

% --- Ambienti semplici (senza riquadro) ---
% Il titolo sta in cima al blocco, anche se il contenuto inizia con un riquadro o
% con codice; e non resta mai da solo in fondo a una pagina.
% Uso: \begin{esempio}[Titolo facoltativo] ... \end{esempio}
\newcommand{\newrunin}[2]{%
  \NewDocumentEnvironment{#1}{o}{%
    \par\Needspace{4\baselineskip}\addvspace{\medskipamount}\noindent
    \textbf{#2}\IfValueT{##1}{ (##1)}\textbf{.}\ \ignorespaces}%
  {\par\addvspace{\medskipamount}}}
\newrunin{esempio}{Esempio}
\newrunin{esercizio}{Esercizio}
\NewTColorBox{osservazione}{}{nota,
  before upper={\textbf{Osservazione.}\ }}

% V e F nelle tavole di verità
\newcommand{\V}{\textsf{V}}
\newcommand{\F}{\textsf{F}}

% --- Codice C ---
% Uso: \begin{codice} ... \end{codice}      codice C numerato
%      \begin{comando} ... \end{comando}    comandi da terminale
%      \begin{risultato} ... \end{risultato} output di un programma
%      \lstinline|printf("ciao");|           codice dentro al testo
\lstdefinestyle{c}{
  language=C,
  basicstyle=\ttfamily\small,
  keywordstyle=\color{blue}\bfseries,
  commentstyle=\color{gray}\itshape,
  stringstyle=\color{red!70!black},
  numbers=left,
  numberstyle=\tiny\color{gray},
  numbersep=8pt,
  columns=flexible,
  keepspaces=true,
  breaklines=true,
  showstringspaces=false,
  tabsize=4,
  morekeywords={include, define},
  % lettere accentate nei commenti
  literate={à}{{\`a}}1 {è}{{\`e}}1 {é}{{\'e}}1 {ì}{{\`i}}1
           {ò}{{\`o}}1 {ù}{{\`u}}1 {È}{{\`E}}1 {À}{{\`A}}1
}
\lstset{style=c}
\newtcblisting{codice}{listing only, nota, left=6mm,
  listing options={style=c}}
\newtcblisting{comando}{listing only, nota,
  listing options={basicstyle=\ttfamily\small, numbers=none, language={},
    columns=flexible, keepspaces=true, breaklines=true}}
\newtcblisting{risultato}{listing only, enhanced jigsaw, breakable,
  colback=black!5, colframe=black, boxrule=0.5pt, sharp corners,
  listing options={basicstyle=\ttfamily\small, numbers=none, language={},
    columns=flexible, keepspaces=true, breaklines=true}}

% --- Diagrammi di flusso (simboli standard) ---
\tikzset{
  terminale/.style = {ellipse, draw, minimum width=2.5cm, minimum height=0.9cm},
  io/.style        = {trapezium, trapezium left angle=70, trapezium right angle=110,
                      draw, minimum width=2.5cm, minimum height=0.9cm},
  processo/.style  = {rectangle, draw, minimum width=2.5cm, minimum height=0.9cm},
  decisione/.style = {diamond, draw, aspect=2, minimum width=2.5cm},
  freccia/.style   = {thick, -{Stealth}}
}

\title{Algebra di Boole e algebra di commutazione}
\author{Marco Cerbella}
\date{Lezione 1 -- 21 settembre 2026}

\begin{document}
\maketitle

\section*{Cenno storico}

L'algebra di Boole fu introdotta da George Boole (1815--1864) nell'Ottocento per
dimostrare affermazioni logiche svolgendo calcoli algebrici: per Boole la
logica va associata alla matematica e non alla filosofia. Nel 1936 Claude
Shannon (1916--2001), studiando al MIT un circuito di oltre cento relè
collegato all'analizzatore differenziale di Bush, intuì che l'algebra di Boole
è un modello formale per l'\emph{analisi} e la \emph{sintesi} di tali sistemi:
due interruttori in \emph{serie} (\emph{parallelo}) si descrivono con AND (OR),
mentre il NOT si realizza con il contatto posteriore di un relè. Dalla sua tesi
(1938) nasce l'algebra di commutazione, base delle reti logiche.

\section{Algebra di Boole}

\begin{definizione}
Un'\emph{algebra di Boole} è una struttura algebrica
\[
  \mathcal{B} = \langle A, +, \cdot, ', 0, 1 \rangle
\]
dove:
\begin{itemize}
    \item $A$ è l'insieme degli elementi su cui l'algebra opera (\emph{supporto});
    \item $+$ è l'operatore di disgiunzione (OR);
    \item $\cdot$ è l'operatore di congiunzione (AND);
    \item $'$ è l'operatore unario di complementazione (NOT);
    \item $0$ è l'elemento neutro rispetto a $+$, $1$ è l'elemento neutro
          rispetto a $\cdot$;
\end{itemize}
e valgono i seguenti assiomi (\emph{postulati di Huntington}).
\end{definizione}

\subsection{Postulati di Huntington}

\begin{enumerate}
    \item \textbf{Commutatività}: $\forall a, b \in A$:
          $a + b = b + a$ \quad e \quad $a \cdot b = b \cdot a$.
    \item \textbf{Elementi neutri}: esistono $0$ e $1$ tali che, per ogni
          $a \in A$, $a + 0 = a$ e $a \cdot 1 = a$.
    \item \textbf{Distributività} (di ciascuna operazione rispetto all'altra):
          $\forall a, b, c \in A$
          \begin{align*}
            a + (b \cdot c) &= (a + b) \cdot (a + c)\\
            a \cdot (b + c) &= (a \cdot b) + (a \cdot c)
          \end{align*}
    \item \textbf{Complemento}: per ogni $a \in A$ esiste ed è unico $a'$ tale
          che $a + a' = 1$ e $a \cdot a' = 0$.
\end{enumerate}

\begin{teorema}[Principio di dualità]
Ogni proposizione derivata dagli assiomi di Huntington resta valida scambiando
ovunque tra loro le operazioni $+$ e $\cdot$ e gli elementi neutri $0$ e $1$.
\end{teorema}

\begin{osservazione}
In particolare se $\langle A, +, \cdot, ', 0, 1 \rangle$ è un'algebra di Boole,
lo è anche $\langle A, \cdot, +, ', 1, 0 \rangle$.
\end{osservazione}

\subsection{Esempi}

\begin{esempio}[Algebra delle classi]
Sia $A$ un insieme e $\mathcal{P}(A)$ l'insieme delle sue parti (tutti i
sottoinsiemi di $A$). La struttura
$\langle \mathcal{P}(A), \cup, \cap, \overline{\phantom{x}}, \emptyset, A \rangle$
è un'algebra di Boole:
\begin{itemize}
    \item $\cup$ e $\cap$ sono commutative;
    \item $\emptyset$ è neutro rispetto a $\cup$ e $A$ rispetto a $\cap$:
          per ogni $X \in \mathcal{P}(A)$, $X \cup \emptyset = X$ e
          $X \cap A = X$;
    \item distributività e unicità del complemento si verificano con i
          diagrammi di Venn.
\end{itemize}
È un'algebra di Boole \emph{chiusa}: $\forall X, Y \in \mathcal{P}(A)$,
$X \cap Y$, $X \cup Y$ e $\overline{X} = A \setminus X$ sono sottoinsiemi di $A$,
quindi elementi di $\mathcal{P}(A)$.
\end{esempio}

\begin{teorema}[Rappresentazione di Stone]
Per ogni algebra di Boole finita $\mathcal{A}$ esiste un insieme finito $X$ tale
che $\mathcal{A}$ è isomorfa a
$\langle \mathcal{P}(X), \cup, \cap, \overline{\phantom{x}}, \emptyset, X \rangle$.
\end{teorema}

Conseguenze:
\begin{itemize}
    \item la cardinalità del supporto di un'algebra di Boole finita è sempre una
          potenza di $2$;
    \item per dimostrare che le algebre di Boole finite godono di una proprietà
          basta dimostrarla per l'algebra delle classi (ad esempio con i
          diagrammi di Venn).
\end{itemize}

\begin{esempio}[Altre algebre di Boole]
$\langle \{0,1\}, \min, \max, 1-x, 0, 1 \rangle$ è un'algebra di Boole; così
come quella ottenuta scambiando le operazioni come nell'osservazione sopra.
\end{esempio}

\subsection{Variabili booleane}

Sia $\mathcal{A} = \langle A, +, \cdot, ', 0, 1 \rangle$ un'algebra di Boole. Una
\emph{variabile booleana} è un simbolo che indica un qualsiasi elemento del
supporto $A$; il suo valore è un elemento di $A$. Le variabili si indicano con
le ultime lettere dell'alfabeto ($x, y, \dots$), gli elementi del supporto con
le prime ($a, b, \dots$) e con i simboli $0$ e $1$.

\subsection{Proprietà dell'algebra di Boole}

\begin{proprieta}
Per ogni $x, y, z$ valgono:
\begin{itemize}
    \item \textbf{associatività}: $x + (y + z) = (x + y) + z$ \ e \
          $x \cdot (y \cdot z) = (x \cdot y) \cdot z$;
    \item \textbf{idempotenza}: $x + x = x$ \ e \ $x \cdot x = x$;
    \item \textbf{elemento nullo}: $x + 1 = 1$ \ e \ $x \cdot 0 = 0$;
    \item \textbf{unicità del complemento}: $x'$ è unico;
    \item \textbf{assorbimento}: $x + x \cdot y = x$ \ e \
          $x \cdot (x + y) = x$;
    \item \textbf{semplificazione}: $x + x' \cdot y = x + y$ \ e \
          $x \cdot (x' + y) = x \cdot y$;
    \item \textbf{involuzione}: $(x')' = x$;
    \item \textbf{leggi di De Morgan}: $(x + y)' = x' \cdot y'$ \ e \
          $(x \cdot y)' = x' + y'$;
    \item \textbf{consenso}:
          $x \cdot y + x' \cdot z + y \cdot z = x \cdot y + x' \cdot z$ \ e \
          $(x + y) \cdot (x' + z) \cdot (y + z) = (x + y) \cdot (x' + z)$.
\end{itemize}
\end{proprieta}

Ciascuna si può verificare con i diagrammi di Venn, oppure dimostrare dagli
assiomi.

\begin{esempio}[Dimostrazione dell'idempotenza]
\begin{align*}
x \cdot x &= (x \cdot x) + 0            && \text{elemento neutro}\\
          &= (x \cdot x) + (x \cdot x') && \text{complemento}\\
          &= x \cdot (x + x')           && \text{distributiva}\\
          &= x \cdot 1                  && \text{complemento}\\
          &= x                          && \text{elemento neutro}
\end{align*}
L'idempotenza di $+$ ($x + x = x$) segue per dualità.
\end{esempio}

\section{Algebra di commutazione}

\begin{definizione}[Algebra di commutazione (\emph{switching algebra})]
È l'algebra di Boole il cui supporto è costituito dai soli valori $\{0, 1\}$:
\[
  \mathcal{B} = \langle \{0,1\}, +, \cdot, ', 0, 1 \rangle .
\]
Ogni variabile può assumere solo i valori $0$ e $1$.
\end{definizione}

È l'algebra usata da Shannon per l'analisi dei circuiti. I tre operatori sono
definiti dalle seguenti tabelle di verità.

\begin{center}
\begin{tabular}{c|cc}
$+$ & 0 & 1 \\ \midrule
0 & 0 & 1 \\
1 & 1 & 1
\end{tabular}
\qquad\qquad
\begin{tabular}{c|cc}
$\cdot$ & 0 & 1 \\ \midrule
0 & 0 & 0 \\
1 & 0 & 1
\end{tabular}
\qquad\qquad
\begin{tabular}{c|c}
$x$ & $x'$ \\ \midrule
0 & 1 \\
1 & 0
\end{tabular}
\end{center}

\subsection{OR esclusivo}

L'\emph{OR esclusivo} (XOR) è l'operatore $\oplus$ definito da
\[
  a \oplus b = a' b + a b' ,
\]
cioè vale $1$ quando $a$ e $b$ sono diversi.

\begin{center}
\begin{tabular}{c|cc}
$\oplus$ & 0 & 1 \\ \midrule
0 & 0 & 1 \\
1 & 1 & 0
\end{tabular}
\end{center}

Dalla definizione, applicando le proprietà dell'algebra di Boole, si ricava:
\begin{align}
x \oplus 1 &= x' & x \oplus 0 &= x \\
x \oplus x &= 0 & x \oplus x' &= 1 \\
x' \oplus y &= x \oplus y' = (x \oplus y)'
\end{align}

\section*{Esercizi per casa}

\begin{esercizio}
Usare i diagrammi di Venn per dimostrare le proprietà (associatività,
idempotenza, \dots) dell'algebra booleana su supporto finito.
\end{esercizio}
\begin{esercizio}
Dimostrare che l'algebra di commutazione è un'algebra di Boole, cioè verificare
i postulati di Huntington.
\end{esercizio}
\begin{esercizio}
Dimostrare, per l'algebra di commutazione, associatività, idempotenza, elemento
nullo, unicità del complemento, assorbimento, semplificazione, involuzione, De
Morgan e consenso.
\end{esercizio}

\end{document}
